\section{文献综述与理论假设}

\subsection{供应链韧性的概念演进与决定因素}

供应链韧性的学术概念经历了从静态到动态的演进过程。早期文献将韧性等同于"恢复能力"（Recovery），即供应链在中断后恢复到原始功能的速度\cite{ponomarov2009understanding}。\cite{pettit2010ensuring}提出了更具综合性的框架，将韧性定义为包含"预判—抵抗—响应—恢复—学习"五个阶段的动态能力谱系，每一阶段对应不同的企业资源与组织能力。这一框架的贡献在于将韧性从被动的"止损"概念提升为主动的"能力构建"范式。

在决定因素方面，运营管理文献识别了若干关键前因变量。供应商多元化通过地理分散和来源分散降低单一节点的中断风险；库存冗余在经典的精益生产框架中曾被视作效率损失，但在韧性导向下被重新评价为"战略性冗余"；关系治理通过长期合作与信任机制降低了供应链伙伴在冲击中的协调摩擦。然而，这些研究普遍基于传统工业经济时代的分析范式，对信息化和数字化条件下的新型韧性来源关注有限。\cite{ivanov2021digital}率先将数字孪生概念引入供应链韧性研究，提出了数据驱动的仿真优化框架，标志着数字化与韧性研究的理论交汇起点。

\subsection{数字化转型的经济效应与供应链应用}

数字化转型对企业的经济影响已获得大量实证支持。在生产率维度，数字化投资通过降低信息处理成本、优化要素配置和促进流程标准化，显著提升了企业全要素生产率。在创新维度，大数据分析和AI辅助研发缩短了产品开发周期，增强了企业对市场需求变化的敏感度。\cite{wu2025impact}基于中国制造业企业数据，采用工具变量法识别了供应链数字化对企业绩效的因果效应，发现数字化水平每提升一个标准差，企业全要素生产率提高约3.2\%。\cite{chen2025supply}利用供应链数字化政策冲击构建双重差分模型，进一步验证了数字化通过降低交易成本和增强信息透明度促进企业创新的传导机制。

供应链领域是数字化应用的前沿场景。ERP系统实现了企业内部供应链流程的集成化管理，MES和WMS系统提升了生产和仓储环节的实时可视化水平，工业互联网平台则打通了上下游企业间的信息壁垒。\cite{bai2025exploring}发现，在数字经济环境下，采用可持续发展导向的数字化创新能够显著降低企业的供应链中断风险。\cite{alyasein2026supply}基于动态能力理论，提出供应链数字化通过提升流程整合能力和创新速度来增强数据驱动的灾害免疫力。这些研究逐步将数字化的分析焦点从"效率"转向"韧性"，但多数仍停留在理论推演或案例研究层面，缺乏基于大样本的严格因果识别。

\subsection{数字化转型与供应链韧性的直接实证}

2024年以来，直接将数字化转型作为供应链韧性核心解释变量的实证文献开始密集出现。\cite{li2025digital}率先使用中国上市公司面板数据，采用固定效应模型发现数字化转型对供应链韧性具有显著正向影响，且该效应在制造业和非制造业子样本中均成立。\cite{qi2024enterprise}从财务视角构建了数字化水平指数，发现数字化水平越高的企业在面临外部冲击时表现出更强的供应链韧性，韧性最强的企业群体在疫情冲击期间的库存恢复速度比韧性最弱群体快约40\%。\cite{zhang2025impact}和\cite{yuan2024effects}分别以供应链整合为中介变量，发现了数字化通过增强内部和外部整合间接提升韧性的传导路径，但在样本规模和中介检验方法上存在局限——前者仅基于212家企业的问卷调查数据，后者未进行正式的Sobel检验或Bootstrap检验。

中文文献中，\cite{zhang2024supply}利用供应链创新与应用试点城市这一准自然实验，采用双重差分法识别了供应链数字化对韧性的因果效应，是中国情境下识别策略最为干净的研究之一。\cite{liude2025supply}基于2012—2022年沪深A股上市公司数据，发现供应链数字化通过信息透明度和协同效率两个渠道增强了韧性。\cite{shi2024digital}将数字化聚焦到数字金融子维度，发现数字金融通过促进技术创新和缓解融资约束间接提升供应链韧性。\cite{song2025valuechain}从价值链视角指出，数字化转型对韧性存在上下游异质性，上游供应商和下游客户对数字化的韧性响应可能存在显著差异。

综观既有实证文献，可以识别出一个清晰的研究空白：现有文章的传导机制讨论最多涵盖一到两条中介路径，且多数未经正式的Sobel检验或Bootstrap置信区间验证。尚未有研究在同一框架下同时检验三条不同方向的传导路径，并系统比较各路径的相对贡献。这构成了本文的核心学术动机。

\subsection{理论框架与研究假设}

本文整合动态能力理论、信息不对称理论和组织柔性理论，构建了一个"双固定效应—三路径中介"的分析框架。理论逻辑如下：数字化转型为企业赋予了三种不同类型的动态能力——信息感知能力（通过数据采集和实时监测）、流程响应能力（通过自动化和智能调度）和战略适应能力（通过数据驱动的创新决策）。这三种能力对应供应链韧性的三个功能维度：感知冲击、响应中断和适应新环境。

据此，我们提出以下递进式的研究假设：

\textbf{H1（主效应假设）}：在其他条件不变的情况下，企业数字化转型水平越高，其供应链韧性越强。

\textbf{H2（中介效应假设）}：数字化转型通过以下三条路径间接增强供应链韧性——\textbf{H2a}：通过提升流程优化度增强供应链响应速度；\textbf{H2b}：通过增强企业创新能力提升供应链适应能力；\textbf{H2c}：通过促进信息共享水平改善供应链预警机制。

\textbf{H3（异质性假设）}：数字化转型对供应链韧性的影响在以下三个维度上存在异质性——企业规模（大型企业的数字化基础设施更完善，但中小企业的边际提升空间可能更大）、财务杠杆（高杠杆企业面临更强的财务约束，数字化可能发挥替代性的风险管理功能）和时间窗口（系统性冲击期间，数字化的保护价值更为突出）。